%=====================================================================
%  DATABASE SYSTEMS --- FALL 2026
%  One-page weekly lecture schedule
%
%  SOURCE OF RECORD: parts/G_schedule.tex in the handout.  Change the
%  handout first, then this page -- the slide pipeline parses the
%  appendix, not this file.
%
%  Follows HEC 2023 (BS Curriculum, Computing Disciplines, p. 58), where
%  the course is Computing Core at semester 2, 4 (3-3) -- the placement
%  the department uses. Its content is HEC 2017's (p. 88) word for word;
%  HEC 2025 lists the course for BSCS and ADCS only and publishes no
%  contents. The department teaches it as CS-103, BSIT, 2nd semester.
%
%  The preamble is shared with the sibling courses' schedule pages.
%
%  Build:  pdflatex DBS_Fall2026_Lecture_Schedule.tex
%=====================================================================
\documentclass[10pt,a4paper]{article}
\usepackage[a4paper,left=1.2cm,right=1.2cm,top=0.7cm,bottom=0.45cm]{geometry}
\usepackage[T1]{fontenc}
\usepackage{lmodern}
\usepackage{microtype}
\usepackage[table]{xcolor}
\usepackage{array,ragged2e}
\usepackage{graphicx}
\usepackage{fontawesome5}
\pagestyle{empty}
\pdfinfo{/Title (Database Systems, Fall 2026: Weekly Lecture Schedule)}
\setlength{\parindent}{0pt}
\setlength{\parskip}{0pt}

%--- palette: shared BSIT palette, part colours darkened a little so
%    white chip text stays legible at 6 pt
\definecolor{primary}{HTML}{1A5276}
\definecolor{secondary}{HTML}{2E86C1}
\definecolor{accent}{HTML}{C0392B}
\definecolor{lightbg}{HTML}{EBF5FB}
\definecolor{neutral}{HTML}{707B7C}
\definecolor{sub}{HTML}{465152}
\definecolor{p1}{HTML}{1A5276}
\definecolor{p2}{HTML}{2874A6}
\definecolor{p3}{HTML}{1E8449}
\definecolor{p4}{HTML}{8E44AD}
\definecolor{p5}{HTML}{117A65}
\definecolor{p6}{HTML}{8A6D0A}

%--- building blocks -------------------------------------------------
\newcommand{\chip}[2]{{\setlength{\fboxsep}{1.3pt}\colorbox{#1}{\makebox[1.45em][c]{%
  \fontsize{6.4}{7}\selectfont\bfseries\color{white}#2}}}}
\newcommand{\chipw}[2]{{\setlength{\fboxsep}{1.3pt}\colorbox{#1}{%
  \fontsize{6.4}{7}\selectfont\bfseries\color{white}\,#2\,}}}
% \ch{part}{number}{title}{what it covers}  -- one chapter, own paragraph
\newcommand{\ch}[4]{\chip{p#1}{#2}\,\textbf{#3}\ {\color{neutral}--}\ {\fontsize{6.9}{8.0}\selectfont\color{sub}#4}\par}
\newcommand{\ex}[1]{{\bfseries\color{accent}#1}}      % exams and the big deliverables
\newcommand{\ev}[1]{{\bfseries\color{primary}#1}}     % quizzes, assignments, project points
\newcommand{\none}{{\color{neutral}---}}
\newcommand{\hdr}[1]{{\bfseries\color{white}#1}}
\newcommand{\lab}[1]{{\color{sub}#1}}

\newlength{\topicw}   % set inside the document, once \linewidth and \tabcolsep are final
\newcolumntype{K}{>{\centering\arraybackslash}p{0.75cm}}
\newcolumntype{T}{>{\RaggedRight\arraybackslash}p{\topicw}}
\newcolumntype{L}[1]{>{\RaggedRight\arraybackslash}p{#1}}
% \leavevmode: a cell that opens with \color puts a whatsit before its paragraph, and
% \vtop then aligns it by its top edge -- one strut lower than a cell that opens with text
\newcommand{\wk}[4]{\leavevmode{\large\bfseries\color{primary}#1} & \leavevmode #2 &
  \leavevmode #3 & \leavevmode #4 \\ \hline}

\newcommand{\fact}[3]{{\setlength{\fboxsep}{4pt}\colorbox{lightbg}{%
  \parbox[t]{\dimexpr#1-8pt\relax}{\RaggedRight\fontsize{6.2}{7}\selectfont\bfseries\color{neutral}\MakeUppercase{#2}\\[1.5pt]%
  \fontsize{7.8}{9.2}\selectfont\bfseries\color{primary}#3}}}}
\newcommand{\tile}[3]{{\colorbox{lightbg}{\parbox[c][0.9cm][c]{\dimexpr(\linewidth-3pt*\numexpr#1-1\relax)/#1-6pt-0.5pt\relax}{%
  \centering{\large\bfseries\color{primary}#2}\\[-0.5pt]{\fontsize{6.6}{7.4}\selectfont\color{sub}#3}}}}}

% \legend{I}{II}{III}{IV}{V}{VI} -- the six part names, as a 3 x 2 grid of colour chips
\newcommand{\legend}[6]{{\fontsize{6.6}{8}\selectfont\color{sub}%
  \begin{tabular}{@{}p{0.31\linewidth}p{0.31\linewidth}p{0.31\linewidth}@{}}
  \chipw{p1}{I}\ #1 & \chipw{p2}{II}\ #2 & \chipw{p3}{III}\ #3 \\[1pt]
  \chipw{p4}{IV}\ #4 & \chipw{p5}{V}\ #5 & \chipw{p6}{VI}\ #6
  \end{tabular}}}

\begin{document}
\fontsize{7.9}{9.3}\selectfont

%--- header ----------------------------------------------------------
\begin{minipage}[c]{1.5cm}%
  \IfFileExists{assets/iub-logo.png}{\includegraphics[height=1.15cm]{assets/iub-logo.png}}{}%
\end{minipage}\hfill
\begin{minipage}[c]{0.65\linewidth}\centering
  {\footnotesize\bfseries\color{primary}THE ISLAMIA UNIVERSITY OF BAHAWALPUR}\\[1pt]
  {\scriptsize\color{neutral}Department of Information Technology}
\end{minipage}\hfill
\begin{minipage}[c]{1.5cm}\raggedleft
  \IfFileExists{assets/dit-logo.png}{\includegraphics[height=1.15cm]{assets/dit-logo.png}}{}%
\end{minipage}

\vspace{3pt}
{\color{secondary}\rule{\linewidth}{1.4pt}}\\[2pt]
{\LARGE\bfseries\color{primary}DATABASE SYSTEMS}\hfill
{\normalsize\bfseries\color{secondary}Weekly Lecture Schedule\ \textbullet\ Fall 2026}\\[1pt]
{\footnotesize\color{sub}Bachelor of Science in Information Technology\ \textbullet\
 \emph{Designing, Querying and Trusting a Relational Database}}\\[1pt]
{\footnotesize\color{sub}CS-103, 2nd semester. One database, \texttt{university.db}, is created in Week 2 and every week afterwards adds to it; by Week 16 it is the project.}\\[1pt]
{\color{secondary}\rule{\linewidth}{0.6pt}}

\vspace{4pt}
\fact{0.235\linewidth}{Credit hours}{4 (3+1): 3 lecture hours and a credited laboratory}\hfill
\fact{0.235\linewidth}{Duration}{16 teaching weeks; midterm Week 8, project due Week 16}\hfill
\fact{0.235\linewidth}{Prerequisite}{Programming Fundamentals; no Data Structures assumed}\hfill
\fact{0.235\linewidth}{Contents}{23 chapters in 6 parts; 11 measured experiments on SQLite}

\vspace{5pt}
%--- the schedule ----------------------------------------------------
\begingroup
\renewcommand{\arraystretch}{0.96}
\setlength{\tabcolsep}{3.5pt}
\setlength{\topicw}{\dimexpr\linewidth-0.75cm-3.55cm-2.85cm-6\tabcolsep\relax}
\arrayrulecolor{primary!18}
\rowcolors{2}{white}{lightbg}
\noindent\begin{tabular}{@{}K T L{3.55cm} L{2.85cm}@{}}
\rowcolor{primary}
\hdr{\faIcon{calendar-alt}} & \hdr{\faIcon{book-open}\ \ Lecture topics} &
\hdr{\faIcon{flask}\ \ Laboratory} & \hdr{\faIcon{clipboard-check}\ \ Assessment} \\
\wk{1}
 {\ch{1}{1}{Why Not Just Files?}{the six problems of a file-based system; database and DBMS}
  \ch{1}{2}{Three Levels, and Why They Are Separate}{external, conceptual, internal; schema and instance; data independence}}
 {\lab{Break a file-based system on purpose; find the three levels in SQLite}}{\none}
\wk{2}
 {\ch{1}{3}{What Is Inside a DBMS, and Who Does What}{one SELECT through six components; the data dictionary; four roles}
  \ch{1}{4}{Your First Database, in One Command}{three statements and a file; a table is 8{,}192 bytes; growth in whole pages}}
 {\lab{Install nothing; create a database and measure what it cost}}{\none}
\wk{3}
 {\ch{2}{5}{The Entity-Relationship Model}{entities, the four attribute distinctions, keys, degree; ER before normalisation}}
 {\lab{Entities and attributes for the university, on paper}}{\none}
\wk{4}
 {\ch{2}{6}{Cardinality, Participation and Weak Entities}{four answers per relationship; (min,max); the section as a weak entity}}
 {\lab{The university ER diagram, finished and argued over}}{\none}
\wk{5}
 {\ch{2}{7}{The Relational Model}{relations, tuples, domains; why a SQL table is not a relation; null}
  \ch{2}{8}{Keys and Integrity Constraints}{entity and referential integrity; foreign keys are off by default}}
 {\lab{Orphan rows in a correctly declared schema}}{\ev{Quiz 1} (Ch 1--6)}
\wk{6}
 {\ch{3}{9}{Mapping an ER Diagram to Tables}{seven rules, applied in order; the M:N table; weak entities and recursion}}
 {\lab{Your diagram becomes eight tables, by rule}}{\ev{Assignment 1 set}}
\wk{7}
 {\ch{3}{10}{SQL: Creating the Schema}{constraints; type affinity against STRICT; an average 20\% wrong}
  \ch{3}{11}{Changing a Schema That Is Already Full}{ALTER TABLE's four forms; the twelve-step rebuild; migrations}}
 {\lab{CREATE TABLE; then change one that is already full}}{\none}
\wk{8}
 {\ch{4}{12}{Relational Algebra}{selection, projection, product, set operations; why ``never'' is a difference}}
 {\lab{Midterm, then algebra by hand on small relations}}{\ex{Midterm} (Ch 1--11)}
\wk{9}
 {\ch{4}{13}{Joins: the Algebra and the SQL}{inner and outer; why NATURAL JOIN returned nothing; automatic indexes}
  \ch{4}{14}{Selecting Rows, and the Trouble with NULL}{three-valued logic; NOT IN returning nothing}}
 {\lab{Lose 73 sections and 19{,}992 rows to joins and NULL, then find them all}}{\ev{Assignment 1 due}}
\wk{10}
 {\ch{4}{15}{Sub-queries, EXISTS and Set Operations}{semi- and anti-joins; correlation was never the cost}
  \ch{4}{16}{Grouping, Aggregation and Window Functions}{what aggregates skip; WHERE against HAVING; RANK}}
 {\lab{One question seven ways; an average moved by a third}}{\ev{Project proposal}}
\wk{11}
 {\ch{5}{17}{Functional Dependencies}{Armstrong's axioms; attribute closure; finding keys; minimal cover}}
 {\lab{The wide report table: its dependencies, closures and key}}{\none}
\wk{12}
 {\ch{5}{18}{Normal Forms: 1NF to BCNF}{the three anomalies, measured; decomposing to BCNF, and in which order}}
 {\lab{Break the running database, then normalise it and measure both}}{\ev{Quiz 2} (Ch 12--17)}
\wk{13}
 {\ch{5}{19}{Decomposition, and Denormalising on Purpose}{lossless joins; dependency preservation; a summary table priced}
  \ch{5}{20}{Views, Triggers and Keeping Data Honest}{a view costs nothing; a trigger costs what its body costs}}
 {\lab{A lossy split that invents facts; a summary table caught lying}}{\ev{Assignment 2 set}}
\wk{14}
 {\ch{6}{21}{Indexes, and Where They Stop Helping}{the B-tree without Data Structures; the crossover; an optimiser that could not tell}}
 {\lab{One million rows, with and without an index}}{\none}
\wk{15}
 {\ch{6}{22}{Transactions, Concurrency and Recovery}{ACID; a lost update; a deadlock; what a commit costs}}
 {\lab{Two writers, one database, and a lost update}}{\ev{Assignment 2 due}}
\wk{16}
 {\ch{6}{23}{Security, NoSQL, and the Semester Project}{SQL injection on your own file; GRANT by contrast; documents against tables}}
 {\lab{Project presentations}}{\ex{Project due; Final} (all)}
\end{tabular}
\endgroup

\vspace{4pt}
%--- assessment weights ----------------------------------------------
{\fontsize{6.6}{8}\selectfont\bfseries\color{neutral}\MakeUppercase{Assessment weights}}\\[2pt]
\tile{6}{10\%}{2 quizzes, Weeks 5 and 12}\hfill
\tile{6}{10\%}{2 assignments}\hfill
\tile{6}{15\%}{Laboratory, weekly from Week 2}\hfill
\tile{6}{20\%}{Mid-term, Week 8}\hfill
\tile{6}{20\%}{Project, built from Week 2}\hfill
\tile{6}{25\%}{Final-term, all chapters}

\vspace{3pt}
%--- legend and notes ------------------------------------------------
\legend{What a Database Is}{Designing It}{From Diagram to Schema}
       {Asking Questions}{Getting the Design Right}{Making It Work, and What Comes Next}

\vspace{1pt}
{\fontsize{6.3}{7.6}\selectfont\color{neutral}%
Follows HEC 2023 (\emph{BS Curriculum, Computing Disciplines}, p.\ 58), where the course is Computing Core at semester 2 with 4 credit hours --- the placement the department uses; HEC 2017 (p.\ 88) placed it at BSIT semester V with the same content word for word, and HEC 2025 lists it for BSCS and ADCS only, with no contents. All 32 distinct clauses of the outline are traced to chapters in the handout's Course Specification, which also records six defects: the text frozen since 2017, the ER model listed after normalisation (the handout reverses them), a duplicated clause, a prerequisite given as ``None'' and as ``Data Structures'' in the same booklet, a mangled CLO table, and a Bloom level that is wrong. HEC names no DBMS in any edition, so SQLite --- already installed with Python, zero setup --- is within the curriculum; what it cannot do (users and GRANT, stored procedures, row locking) is taught by contrast. Chapter numbers refer to the handout, \emph{Database Systems --- Designing, Querying and Trusting a Relational Database} (213 pp., 23 chapters, 7 appendices); Appendix~G is the source of record for this page. Eleven chapters carry a measured experiment whose figures were produced on the machine described in Appendix~A; the sources are in \texttt{code/}, \texttt{code/run\_all.py} re-checks every claim, and each is reproducible. The security laboratory is confined to the student's own database file; testing any other system is an offence under PECA 2016.\par}

\end{document}
